\section{Plan de Trabajo}

El trabajo se organiza en cuatro fases. Las dos primeras preparan los datos y
las variables; las dos últimas construyen, ajustan y comparan los modelos. Algunas
actividades se solapan intencionalmente: el análisis exploratorio arranca en la
fase~2 porque los hallazgos descriptivos orientan qué variables conviene incluir,
y la redacción del informe comienza en la fase~4 antes de que termine el análisis
para no acumular toda la escritura al final.

\begin{table}[H]
  \centering
  \caption{Plan de trabajo por fases y actividades}
  \label{tab:plan}
  \renewcommand{\arraystretch}{1.35}
  \begin{tabularx}{\linewidth}{c l X c}
    \toprule
    \textbf{Fase} & \textbf{Actividad} & \textbf{Descripción / Herramientas} & \textbf{Semanas} \\
    \midrule
    \multirow{2}{*}{1}
      & Recolección y limpieza
      & Descarga de registros SECOP~II vía API o CSV. Normalización de campos,
        imputación de nulos, deduplicación y corrección de codificación territorial.
        Python, Pandas, Requests.
      & 1\,--\,5 \\
      & Variable objetivo
      & Definición del \textit{proxy} de incumplimiento: adiciones contractuales,
        suspensiones reiteradas y desfase entre fecha pactada y fecha real de
        liquidación. Análisis de sensibilidad del umbral elegido.
      & 3\,--\,5 \\
    \midrule
    \multirow{2}{*}{2}
      & Ingeniería de características
      & Selección, transformación y codificación de variables predictoras (anticipos,
        modalidad, proponentes, ubicación). Análisis exploratorio para detectar
        correlaciones y desbalance de clases. Scikit-learn, Feature-engine.
      & 5\,--\,8 \\
      & \textit{Pipeline} de preprocesamiento
      & Construcción del flujo reproducible: imputación, escalado estándar y
        codificación \textit{one-hot}. Serialización con \texttt{joblib} para
        reutilización en inferencia. Scikit-learn Pipeline.
      & 6\,--\,8 \\
    \midrule
    \multirow{2}{*}{3}
      & Entrenamiento de modelos
      & Ajuste de Regresión Logística (línea base), Random Forest y XGBoost.
        Optimización con \texttt{RandomizedSearchCV} ($n\_iter = 50$).
        Google Colab, Scikit-learn, XGBoost.
      & 8\,--\,12 \\
      & Validación y métricas
      & Validación cruzada estratificada $k = 5$. Reporte de AUC-ROC, F1-score,
        precisión, \textit{recall} y curvas PR para cada modelo.
      & 11\,--\,13 \\
    \midrule
    \multirow{2}{*}{4}
      & Análisis comparativo
      & Comparación de modelos. Análisis de importancia con SHAP y
        \textit{permutation importance}. Selección y justificación del modelo final.
      & 13\,--\,15 \\
      & Informe y monografía
      & Redacción del informe técnico y la monografía. Organización del repositorio
        con código, datos de muestra y ambiente reproducible.
      & 14\,--\,16 \\
    \bottomrule
  \end{tabularx}
\end{table}

El repositorio en Git documentará cada experimento con su semilla, versión de
librerías y métricas obtenidas. Eso permite retomar cualquier punto del trabajo
sin tener que repetir el procesamiento completo.